\section{Evaluation Protocol and Placeholder Results}
\label{sec:experiments}

\defname{} is a benchmark, so ``evaluation'' means \emph{validating the instrument} and
\emph{demonstrating the headline finding}. This section specifies the protocol; the
numerical tables and plots are clearly labelled placeholders, and all quantitative
outcomes are marked \emph{(hypothesis)} until produced. \Cref{fig:expsetup} summarizes the
design as a pipeline from datasets and variables to metrics.

\subsection{Research questions}
\begin{itemize}[leftmargin=1.4em]
\item[\textbf{RQ1}] \emph{Construct validity.} Does $\NRPs$ behave as
\Cref{thm:metric} predicts (monotone in $\PNAs$ and $1-\ASRs$, reducing to ASB at $N=1$),
and do alternative aggregators (\Cref{rem:metric}) change configuration rankings?
\item[\textbf{RQ2}] \emph{Evaluator integrity (hypothesis).} How large is the
deterministic-versus-LLM-judge gap---i.e., how often can an injection hijack a judge but
not the deterministic checker (\Cref{thm:integrity,prop:judgehijack})?
\item[\textbf{RQ3}] \emph{Defense degradation (headline hypothesis).} Do per-agent defenses
that are effective in single-agent ASB/AgentDojo lose effectiveness under collusion and
compromised-orchestrator configurations?
\item[\textbf{RQ4}] \emph{Topology and propagation (hypothesis).} Does the measured $\CPR$
follow the topological ordering predicted by \Cref{prop:chain,prop:star,prop:tree,prop:mesh}
(chain $<$ tree$_{\text{sub}}$ $<$ star $<$ mesh), and is the orchestrator the critical node?
\item[\textbf{RQ5}] \emph{Collusion and Sybils (hypothesis).} Is there a measurable
collusion advantage $\Delta_{\mathrm{coll}}$ (\Cref{def:csr,thm:collusion}), and how does
Sybil pressure affect selection/voting/memory outcomes?
\item[\textbf{RQ6}] \emph{Scalability and cost.} How do the metrics and the token/API and
latency costs scale with the agent count $N\in\{3,\dots,50\}$ and \corb{} rounds?
\end{itemize}

\subsection{Models, scale, and tasks}
\textbf{Agent backbones.} Open models---Llama-3~\cite{dubey2024llama3},
Qwen2.5~\cite{qwen2025qwen25}, and Mistral~\cite{jiang2023mistral7b}---and proprietary
models---GPT-4o~\cite{openai2024gpt4o} and Claude-3.5---so that findings are not tied to a
single family. \textbf{Scale/topology.} $3$--$50$ agents across star, chain, tree, and mesh.
\textbf{Tasks.} Forced-coordination tasks (\Cref{def:forced}) instrumented with registered
$(\Predutil,\Predsec)$ pairs; benign difficulty tuned so $\PNAs$ is neither saturated nor
zero, with benign task pools drawn in the spirit of agent benchmarks such as
AgentBench~\cite{liu2024agentbench}, GAIA~\cite{mialon2024gaia},
WebArena~\cite{zhou2024webarena}, and SWE-bench~\cite{jimenez2024swebench}, and reasoning
pools such as MMLU~\cite{hendrycks2021measuring} and GSM8K~\cite{cobbe2021gsm8k} for worker
subtasks.

\subsection{Adversary sweeps}
We sweep each knob of $\config_{\mathrm{adv}}$ (\Cref{eq:adv-config}) while holding the
others fixed: adversary fraction $\betafrac\in\{0,0.1,\dots,0.5\}$; independent versus
colluding; static versus adaptive; orchestrator trusted versus compromised; malicious tools
on/off; Sybil count $\in\{0,1,2,4\}$; and poisoned-memory on/off. Red-team primitives are
instantiated from the single-agent literature---indirect prompt
injection~\cite{greshake2023indirect,zhan2024injecagent}, instruction
override~\cite{perez2022ignore}, universal triggers~\cite{zou2023universal}, and
memory poisoning~\cite{chen2024agentpoison}---and composed across roles by the harness.

\subsection{Baselines and defenses}
As stated in the source design, baselines are: (i) \textbf{ASB re-run in single-agent
mode}~\cite{zhang2025asb}; (ii) \textbf{AgentDojo}~\cite{debenedetti2024agentdojo};
(iii) \textbf{no-defense}; and (iv) \textbf{per-agent defenses lifted into the multi-agent
setting}---delimiting/sandwiching, paraphrase, a prompt-injection detector, instructional
prevention, and guardrails such as Llama Guard~\cite{inan2023llamaguard} and NeMo
Guardrails~\cite{rebedea2023nemo}, plus TrustAgent-style safety
prompting~\cite{hua2024trustagent}. Where relevant we also include multi-agent defenses as
reference points (AutoDefense~\cite{zeng2024autodefense}, G-Safeguard~\cite{yu2025gsafeguard})
and an oracle upper bound that scores with ground-truth compromise labels.

\subsection{Metrics and ablations}
Metrics are the full suite of \Cref{sec:metric-suite}: $\ASRs$, $\PNAs$, $\NRPs$, $\CSR$,
$\Delta_{\mathrm{coll}}$, $\CPR$, coordination quality, time-to-detection, recovery rate,
token/API and latency cost, scalability, and human-intervention rate. The registered
\textbf{ablations} are: collusion on/off; orchestrator trusted vs.\ compromised; emulated
vs.\ real tools (bounding the $\kappaem\approx0.48$ realism cost); and deterministic vs.\
LLM-judge evaluation (quantifying evaluator-hijack risk, RQ2). Each metric is estimated with
$M$ episodes chosen by \Cref{thm:sample} for a target $(\epsilon,\alpha)$, with Benjamini--Hochberg
control across the sweep grid.

\subsection{Implementation and hyperparameters}
The harness is model-agnostic (any chat/completions backbone behind a common interface),
runs episodes with horizon $H$ per task, and logs every message, tool call, memory
operation, and state transition into $\Trace$. Default study parameters (to be finalized
with compute budget): $M{=}200$ episodes per configuration at target $\epsilon{=}0.05$,
$\alpha{=}0.05$; $R\!\le\!20$ \corb{} rounds; temperature and decoding fixed per backbone
and reported; random seeds logged. All defenses and attacks are versioned, and every reported
number is a row of the \corb{} trajectory $\mathcal{L}$ (\Cref{alg:corb}).

\subsection{Placeholder result tables}
\Cref{tab:degradation,tab:topology} give the \emph{structure} of the two headline results.
Their cells are placeholders.

\begin{table}[t]
\centering
\caption{\textbf{Placeholder / expected trend only; replace with actual experimental results
before submission.} Headline degradation study (RQ3): system-level $\ASRs\downarrow$ /
$\NRPs\uparrow$ are better. Each defense is measured single-agent (ASB/AgentDojo regime) and
then under multi-agent collusion and compromised-orchestrator configurations. Expected
\emph{(hypothesis)} trend: a defense's advantage over no-defense shrinks from left to right.}
\label{tab:degradation}
\small
\begin{tabular}{@{}lcccccc@{}}
\toprule
& \multicolumn{2}{c}{Single-agent} & \multicolumn{2}{c}{MA + collusion} & \multicolumn{2}{c}{MA + compromised orch.} \\
\cmidrule(lr){2-3}\cmidrule(lr){4-5}\cmidrule(lr){6-7}
Defense & $\ASRs$ & $\NRPs$ & $\ASRs$ & $\NRPs$ & $\ASRs$ & $\NRPs$ \\
\midrule
No-defense              & --- & --- & --- & --- & --- & --- \\
Delimiting/sandwich     & --- & --- & --- & --- & --- & --- \\
Paraphrase              & --- & --- & --- & --- & --- & --- \\
PI detector             & --- & --- & --- & --- & --- & --- \\
Instructional prevention& --- & --- & --- & --- & --- & --- \\
Guardrails              & --- & --- & --- & --- & --- & --- \\
\midrule
\emph{Oracle (upper bound)} & --- & --- & --- & --- & --- & --- \\
\bottomrule
\end{tabular}
\end{table}

\begin{table}[t]
\centering
\caption{\textbf{Placeholder / expected trend only; replace with actual experimental results
before submission.} Topology study (RQ4): measured compromise-propagation rate $\CPR$ and
$\ASRs$ by topology at fixed $\betafrac$. Expected \emph{(hypothesis)} ordering, from
\Cref{prop:chain,prop:star,prop:tree,prop:mesh}: chain $<$ subcritical tree $<$ star $<$
mesh, with the orchestrator (star hub) as the critical node.}
\label{tab:topology}
\small
\begin{tabular}{@{}lccccc@{}}
\toprule
Topology & $\CPR$ & $\ASRs$ & $\NRPs$ & Time-to-detect & Recovery \\
\midrule
Chain & --- & --- & --- & --- & --- \\
Tree (subcritical, $b\pinf<1$) & --- & --- & --- & --- & --- \\
Tree (supercritical, $b\pinf>1$) & --- & --- & --- & --- & --- \\
Star & --- & --- & --- & --- & --- \\
Mesh & --- & --- & --- & --- & --- \\
\bottomrule
\end{tabular}
\end{table}

\Cref{fig:results} plots the \emph{conceptual expected trends} for RQ3--RQ6 (defense
degradation with collusion, $\CPR$ versus topology, collusion advantage versus $\betafrac$,
and cost versus $N$). They are drawn from the theory of \Cref{sec:theory} and are labelled
as conceptual, not measured.

\subsection{Reproducibility checklist}
We will release: the environment harness and orchestration layer; the task suite with
registered $(\Predutil,\Predsec)$ predicates; the adversary knobs and red-team primitives;
the deterministic checker and the optional off-critical-path judge; the \corb{} driver; all
prompts, decoding settings, seeds, and model/version identifiers; and the full \corb{}
trajectory logs $\mathcal{L}$ underlying every reported number. Because scoring is
deterministic given a trace, results are reproducible from the logs even when backbone APIs
change. The emulated-versus-real-tools and deterministic-versus-judge ablations are shipped
as first-class configurations so that reviewers can reproduce the validity checks, not only
the headline numbers.
